\chapter*{Dunkeld}\stepcounter{chapter}\phantomsection\addcontentsline{toc}{chapter}{Dunkeld}

\DndDropCapLine{D}{\entryfont unkeld has been utterly destroyed. Whatever struck the hamlet did so recently and with overwhelming violence; scattered embers still glow among collapsed homes, and thin trails of smoke drift above the ruins. Corpses remain where inhabitants were cut down, several deliberately mutilated or defaced after death. There are no survivors to be heard, nor any obvious signs that someone has returned.}

\begin{DndReadAloud}
	The road emerges from the trees, and the remains of Dunkeld stretch out before you.

	Blackened timbers jut from collapsed houses. Here and there, orange embers still smoulder beneath heaps of charred wood, sending thin wisps of smoke into the grey air. Bodies lie abandoned between the ruins, some bearing wounds far beyond those that killed them. Nothing moves among them.

	Ahead, the road widens into what must once have been the village plaza. A great stone figure lies shattered across the cobbles, torn from the pedestal upon which it once stood. Its face has been hacked away and its form mutilated beyond recognition.

	The pedestal, however, remains.

	Through soot and scratches, its inscription is still legible:

	\textit{"Dunkeld's Protector: Ewan MacRae."}
\end{DndReadAloud}

\section*{Search the Ruins}\phantomsection\addcontentsline{toc}{section}{Search the Ruins}
{\entryfont The characters can spend time searching the destroyed buildings of Dunkeld for anything that survived the attack. Although most valuables have been destroyed, buried beneath rubble, or lost to the fires, a careful search may uncover useful equipment and personal belongings.

Choose one character to lead the search. That character makes an Intelligence (Investigation) Check. Another character may use the Help Action, granting Advantage on the check.

Use the result to determine how many times the party may roll on the \hyperref[table:DunkeldLoot]{\LinkFont{Dunkeld Loot Table}} and whether they receive a bonus to those rolls.}

\begin{DndTable}[header=Search the Ruins]{llX}
	Investigation Check	& Loot Rolls	& Bonus to d100 Rolls	\\
	9 or lower			& 1				& -						\\
	10 - 12				& 2				& -						\\
	13 - 15				& 2				& +5					\\
	16 - 18				& 3				& +5					\\
	19 - 21				& 3				& +10					\\
	22+					& 4				& +10					\\
\end{DndTable}

\subsection*{Rolling for Loot}
{\entryfont For each Loot Roll, roll a d100 and add the bonus granted by the Investigation check. A result cannot exceed 100.

The corresponding entry on the Dunkeld Loot Table is found somewhere among the ruins. The DM can decide exactly where the item is discovered or describe an appropriate location based on the rolled result.}

\begin{DndTable}[header=Dunkeld Loot Table]{lX}\label{table:DunkeldLoot}
	d100	& Loot	\\
	01 - 04	& \textbf{Charred Family Portrait.} A small painted portrait of a smiling Dunkeld family. The edges are badly burned.\\
	05 - 09	& \textbf{Child's Wooden Knight.} A hand-carved toy knight with a faded blue cloak. One arm has broken off.\\
	10 - 12	& \textbf{Broken Silver Locket.} A battered locket containing two tiny portraits. The silver is worth 2 GP.\\
	13 - 15	& \textbf{Dunkeld Drinking Horn.} Decorated with simple carvings of trees and the River Tay worth 2 GP.\\
	16 - 18	& \textbf{Scattered Coin Purse.} Beneath a collapsed cupboard lie 2d6 SP and 1d4 GP.\\
	19 - 21	& \textbf{Carved Pipe.} A well-used wooden smoking pipe decorated with tiny running deer worth 3 GP.\\
	22 - 27	& \textbf{Copper Bracelet.} Found beneath the remains of a collapsed bed worth 5 GP.\\
	28 - 30	& \textbf{Hunter's Trophy.} A polished antler-handled knife. It functions as a normal Dagger and is worth 5 GP.\\
	31 - 36	& \textbf{Surviving Provisions.} A sealed storage container holds 1d4 days of Rations.\\
	37 - 42	& \textbf{Coil of Rope.} 50 feet of Hempen Rope, slightly blackened but perfectly usable.\\
	43 - 51	& \textbf{Hunting Supplies.} A leather satchel containing 20 Arrows and 1d4 days of Rations.\\
	52 - 57	& \textbf{Hunter's Trap.} An intact Hunting Trap, still wrapped for transport.\\
	58 - 63	& \textbf{Herbalist's Supplies.} A surviving Herbalism Kit together with several bundles of dried local herbs.\\
	64 - 66	& \textbf{Locked Strongbox.} The small iron box can be opened with a DC 12 Dexterity check using Thieves' Tools or broken open. It contains 3d6 GP.\\
	67 - 69	& \textbf{Merchant's Purse.} Hidden beneath a loose floorboard are 2d10 GP and a small piece of amber worth 10 GP.\\
	70 - 75	& \textbf{Fine Silverware.} A wrapped set of surviving silver utensils bearing a Dunkeld family crest worth 20 GP.\\
	76 - 78	& \textbf{Silvered Ammunition.} A hunter's locked cabinet contains 1d6 Silvered Arrows.\\
	79 - 81	& \textbf{Antitoxin.} A carefully sealed vial of Antitoxin survives inside an apothecary's padded case.\\
	82 - 87	& \textbf{Potion of Healing.} A single Potion of Healing lies inside a shattered medicine cabinet, somehow still intact.\\
	88 - 92	& \textbf{Dunkeld Hunter's Bow.} A beautifully maintained Longbow. It is nonmagical but worth 75 GP.\\
	93 - 96	& \textbf{Candle of the Deep.} A strange blue candle stored inside a watertight metal case.\\
	97 - 98	& \textbf{Potion Cache.} A concealed compartment in an apothecary contains 2 Potions of Healing and 1 Antitoxin.\\
	99 - 100& \textbf{Dunkeld Survivor's Cache.} A deliberately concealed emergency cache containing 30 gp, 2 Potions of Healing, and 2 Vials of Holy Water.
\end{DndTable}

\begingroup
	\DndSetThemeColor[PhbTan]
	\begin{DndComment}{DM Annotation}
		Only one Investigation check should be made for the entire hamlet. The check represents the combined effort of the party rather than a search by a single character.

		The loot found here is intentionally minor. Dunkeld has already been devastated, and searching its ruins should primarily reward exploration with useful supplies, small valuables, and occasional unusual finds rather than major treasure.
	\end{DndComment}
\endgroup

\section*{Dark Forest of Tay}\phantomsection\addcontentsline{toc}{section}{Dark Forest of Tay}
{\entryfont Beyond Dunkeld, the road quickly gives way to the dense and gloomy Dark Forest of Tay. The woods are eerily quiet, with many of the old hunting trails now abandoned.

As the party ventures deeper, they begin to find signs of recent violence: discarded equipment, disturbed undergrowth, and the mauled remains of hunters.}

\begin{DndReadAloud}
	The clearing is eerily quiet.

	Four bodies lie where they fell among churned earth and shattered saplings. Broken bows, snapped spear shafts, and torn packs are scattered throughout the area.

	The corpses bear terrible wounds. Armor has been crushed inward, ribs splintered, and deep punctures mark several bodies.

	Most unsettling of all, the hunters seem to have died facing different directions, as though attacked from several sides at once.
\end{DndReadAloud}
\subsection*{Clues}\phantomsection\addcontentsline{toc}{subsection}{Clues}
{\entryfont The following clues can be discovered while investigating the remains of the hunters and the surrounding area. Some are immediately apparent, while others require a successful ability check. The clues should gradually allow the party to reconstruct what happened without revealing the answer outright.}

\subsubsection*{Clue 1: The Bloodied Journal}
{\entryfont Near one of the bodies lies a \hyperref[resource:BloodiedHuntersJournal]{\LinkFont{small, leather-bound journal}}, partially soaked in blood but still readable.

The journal contains notes from the hunters' final days in the Dark Forest of Tay and can be given directly to the players.

\subparagraph{Reveals} The hunters entered the forest pursuing what they believed to be a single dangerous creature. They were uncertain whether their quarry was a living beast or something undead.}

\subsubsection*{Clue 2: Strange Wounds}
{\entryfont The hunters' bodies are badly mauled. Bones have been broken, armor crushed inward, and several bodies bear unusually deep puncture wounds.

\paragraph{DC 12 Medicine} Many of the injuries were caused by tremendous blunt force, consistent with being struck or trampled by a large creature.

\paragraph{DC 16 Medicine} The deepest wounds are narrow punctures driven with tremendous force. They were likely caused by a large horn.}

\subsubsection*{Clue 3: The Undergrowth}
{\entryfont The vegetation surrounding the clearing has been heavily disturbed. Bushes have been flattened, branches snapped, and patches of undergrowth torn apart.

\paragraph{DC 12 Nature} The damage was caused suddenly and with considerable force. Several bushes were struck hard enough to snap mature branches rather than merely being pushed aside.

\paragraph{DC 16 Nature} Several plants bear narrow punctures through stems and branches at roughly chest height. The damage is consistent along several different paths into the clearing.}

\subsubsection*{Clue 4: The Battlefield}
{\entryfont The clearing itself is covered in footprints, broken equipment, and signs of a desperate struggle.

\paragraph{DC 14 Survival} The hunters initially positioned themselves against a single approaching creature.

\paragraph{DC 16 Survival} More tracks entered the clearing after the fighting had already begun, approaching from several directions.

\paragraph{DC 18 Survival} At least three large creatures took part in the final moments of the battle. The hunters were surrounded.}

\subsubsection*{Clue 5: Holy Equipment}
{\entryfont Among the hunters' ordinary equipment are several unusual items:
\begin{itemize}
	\item Silvered spearheads
	\item Blessed charms
	\item Vials of holy water
\end{itemize}

\paragraph*{DC 12 Religion} The equipment was deliberately chosen for confronting undead creatures.

\subparagraph{Reveals} Although the hunters did not know exactly what they were pursuing, they considered undeath likely enough to prepare for it.}

\subsection*{Finding the Fifth Hunter}\phantomsection\addcontentsline{toc}{subsection}{Finding the Cache}
{\entryfont If the party has read the Bloodied Journal, they know that five hunters entered the Dark Forest of Tay. Yet only four bodies lie within the clearing.

A \textbf{DC 15 Survival Check} reveals a faint trail leading away from the clearing. Broken twigs, disturbed leaves, and occasional drops of dried blood suggest that someone badly wounded managed to crawl away from the fighting.

Following the trail for a short distance leads to the body of the fifth hunter, concealed beneath the roots and brush of a large tree. He appears to have succumbed to his wounds shortly after escaping the clearing.

Clutched in one hand is a bloodstained note:}
\begin{BloodiedNote}[font=\handwrittenfont, fontsize=18pt]
	If anyone finds this...

	We were wrong.

	There was never just one.

	I hid the gear where they would not find it.

	May the Light guide a steadier hand than mine.
\end{BloodiedNote}
{\noindent\entryfont Nearby, hidden inside a hollow tree \textbf{(DC 13 Intelligence (Investigation) Check)}, rests a leather case containing:
\begin{itemize}
	\item A set of finely crafted hunting gear:
	\begin{itemize}
		\item Ring of Animal Influence
		\item Bridle of Capturing
	\end{itemize}
	\item 2 days of preserved rations
	\item Arrow of Slaying (Beast)
	\item Arrow of Slaying (Undead)
\end{itemize}}